\documentclass[11pt]{article}
\usepackage[margin=1in]{geometry}
\usepackage{amsmath,amssymb,amsthm}
\usepackage[colorlinks=true,urlcolor=blue]{hyperref}
\newtheorem{theorem}{Theorem}
\newtheorem{lemma}{Lemma}
\newtheorem{corollary}{Corollary}
\newcommand{\R}{\mathbb R}
\newcommand{\rank}{\operatorname{rank}}
\newcommand{\Tr}{\operatorname{tr}}
\newcommand{\pdet}{\operatorname{pdet}}
\title{A positive three-column response and a four-vertex\protect\\
obstruction to inherited normalized quadrature cutoffs}
\author{Partitioning frontier research session}
\date{November 15, 2026}
\begin{document}
\maketitle

\paragraph{Status.}
These are proved identities and a counterexample, rather than a new
partitioning algorithm. The matrix-function integral, Schur-complement
update, and operator geometric mean are classical. We record their exact
combination and the explicit obstruction because they decide whether the
repository's fixed-mass fractional update proof transfers to degree
normalization. Historical novelty has not been established.

\section{The normalization-aware form}
For a connected weighted loopless undirected graph, let $L$ be its
Laplacian and $D$ its positive degree matrix. Fix $0<\alpha<1$ and set
\[
 H_\alpha(L,D)=D^{1/2}(D^{-1/2}LD^{-1/2})^\alpha D^{1/2},
 \qquad c_\alpha=\frac{\sin(\pi\alpha)}\pi.
\]
Its kernel is $\R\mathbf1$. Throughout, fractional powers vanish on the
kernel. A positive edge edit has
\[
 L'=L+\delta bb^T,\qquad
 D'=D+\delta(e_ue_u^T+e_ve_v^T),\qquad b=e_u-e_v,\quad\delta>0.
\]
For $t>0$, write
\[
 K_t=L+tD,\qquad
 F_t(L,D)=D-tD K_t^{-1}D
 =\frac1t[L-LK_t^{-1}L].
\]
The scalar fractional-power formula, applied spectrally, gives
\begin{equation}\label{eq:integral}
 H_\alpha(L,D)=c_\alpha\int_0^\infty t^{\alpha-1}F_t(L,D)\,dt.
\end{equation}
Both endpoint integrals converge in ordinary matrix norm, since
$0\preceq F_t\preceq D$ and $0\preceq F_t\preceq L/t$.

The companion note \texttt{higher-cheeger/markov-form-budget.tex}
proves joint monotonicity and the exact relative determinant identity
\begin{align}\label{eq:budget}
 B_\alpha&:=\log\frac{\pdet H_\alpha(L',D')}{\pdet H_\alpha(L,D)}\\
 &=\alpha\log(1+\delta b^TL^\dagger b)
 +(1-\alpha)\log\frac{(1+\delta/d_u)(1+\delta/d_v)}
 {1+2\delta/\Tr D}.\notag
\end{align}
In particular, $\Delta H=H_\alpha(L',D')-H_\alpha(L,D)\succeq0$,
and the number of eigenvalues greater than $\varepsilon$ of
$H_\alpha^{-1/2}\Delta H H_\alpha^{-1/2}$ is strictly smaller than
$B_\alpha/\log(1+\varepsilon)$ when that count is positive.
All inverses in this last statement are on $\mathbf1^\perp$.

\section{An explicit positive resolvent response}
\begin{lemma}[Schur-complement response]
Let $\mathcal A=\left[\begin{smallmatrix} A&B\\B^T&C\end{smallmatrix}\right]$
be symmetric with $C\succ0$, and add
$\left[\begin{smallmatrix} U\\V\end{smallmatrix}\right]
\left[\begin{smallmatrix} U\\V\end{smallmatrix}\right]^T$.
The change in its Schur complement onto the first block is
\[
 (U-BC^{-1}V)(I+V^TC^{-1}V)^{-1}(U-BC^{-1}V)^T.
\]
\end{lemma}
\begin{proof}
Use the invertible change of variables
$y=z-C^{-1}B^Tx$ in the block quadratic form. Its base cross term
vanishes, while the added term becomes
$\|(U-BC^{-1}V)^Tx+V^Tz\|^2$. Minimizing over $z$ gives the display,
by the elementary inverse identity
$I-V^T(C+VV^T)^{-1}V=(I+V^TC^{-1}V)^{-1}$.
This also proves positivity and the rank bound without any nonsingular
assumption on the full block matrix.
\end{proof}

\begin{theorem}[Three-column positive integrand]\label{thm:three}
Define $n$-by-$3$ matrices
\[
 V_t=\sqrt\delta[\,b,\ -\sqrt t\,e_u,\ -\sqrt t\,e_v\,],
\]
\[
 G_t=\sqrt\delta[\,\sqrt t\,D K_t^{-1}b,\
                    LK_t^{-1}e_u,\ LK_t^{-1}e_v\,].
\]
Then the exact one-edge response at every positive shift is
\begin{equation}\label{eq:factor}
 F_t(L',D')-F_t(L,D)
   =G_t(I+V_t^TK_t^{-1}V_t)^{-1}G_t^T.
\end{equation}
It is positive semidefinite, has rank at most three, and annihilates
constants. Consequently $\Delta H$ is the positive integral of these
matrices with weight $c_\alpha t^{\alpha-1}$.
The rank-three assertion concerns each integrand, not the integrated
response.
\end{theorem}
\begin{proof}
The form $F_t$ is the Schur complement of the two-layer network
\[
 \mathcal A_t=\begin{bmatrix}D&-D\\-D&D+L/t\end{bmatrix}.
\]
An edge edit adds three PSD rank-one forms: an internal edge
$\sqrt{\delta/t}[0;b]$, and two terminal-to-internal edges
$\sqrt\delta[e_u;-e_u]$, $\sqrt\delta[e_v;-e_v]$.
Apply the lemma with $B=-D$, $C=K_t/t$, and these columns.
The resulting denominator is $I+V_t^TK_t^{-1}V_t$, and its numerator
columns are exactly $G_t$, since $I-tDK_t^{-1}=LK_t^{-1}$.
For the first column,
$\mathbf1^TDK_t^{-1}b=t^{-1}\mathbf1^Tb=0$; the other columns
are in the range of $L$. This proves the kernel assertion.
\end{proof}

This is a classical Kron-reduction/low-rank-update calculation. It does
provide a concrete interface for a prospective normalized relative
algorithm: three shifted-system columns and one $3$-by-$3$ core per node.
It does not establish a sufficient quadrature range or a charged algorithm.

\section{A four-vertex obstruction}
Consider the weighted path $1$--$2$--$3$--$4$ with edge weights
$(w,h,1)$, where $w>0$ is fixed and $h\downarrow0$. Let
$L_{w,h},D_{w,h}$ be its Laplacian and degree matrix, and set
$x=(1,1,-1,-1)^T$.

\begin{lemma}[Exact slow-scale asymptotic]\label{lem:slow}
For every fixed $w>0$ and $0<\alpha<1$,
\[
 x^TH_\alpha(L_{w,h},D_{w,h})x
 =4h^\alpha\left(\frac{w+1}{2w}\right)^{\alpha-1}(1+o(1)).
\]
\end{lemma}
\begin{proof}
The normalized adjacency has spectrum $\{1,-1,\rho,-\rho\}$, where
\[
 \rho=\sqrt{\frac{w}{(w+h)(1+h)}}.
\]
Indeed its characteristic polynomial is even, it has eigenvalues
$\pm1$ by bipartiteness, and its determinant is the product of the
normalized weights of its first and third edges, squared. Thus the
normalized Laplacian eigenvalues are $0,1-\rho,1+\rho,2$, and
\[
 \lambda_{\rm slow}=1-\rho
 =h s+O(h^2),\qquad s=\frac{w+1}{2w}.
\]
At $h=0$, use the orthonormal component-constant vectors
\[
 q_L=(1,1,0,0)^T/\sqrt2,\qquad
 q_R=(0,0,1,1)^T/\sqrt2.
\]
The limiting stationary vector has coefficients $(\sqrt{2w},\sqrt2)$,
and the limiting slow vector has coefficients
$(\sqrt2,-\sqrt{2w})/\sqrt{2w+2}$. The vector
$D_{w,h}^{1/2}x$ tends to a vector with coefficients
$(\sqrt{2w},-\sqrt2)$. Its squared projection onto the slow vector
therefore tends to $8w/(w+1)=4/s$.
The high spectral subspace is uniformly separated from zero. Its
projection of $D_{w,h}^{1/2}x$ is $O(h)$, since that projection
vanishes at $h=0$ and the corresponding spectral projector is analytic
near zero. Its energy contribution is consequently $O(h^2)$.
The stationary component contributes zero. The remaining contribution
is $(hs+O(h^2))^\alpha(4/s+o(1))$, as asserted.
\end{proof}

Now double the first edge of the path with weights $(1,h,1)$, producing
$(2,h,1)$. Put $H_h=H_\alpha(L_{1,h},D_{1,h})$ and $H'_h$ for the
updated form. The edge is a bridge, so its leverage before doubling is
exactly $\tau=1$. Both endpoint degrees and the graph size stay bounded
above and below, independently of $h$.

\begin{theorem}[Failure of an inherited gap-free low cutoff]\label{thm:barrier}
The relative response on $x$ satisfies
\begin{equation}\label{eq:limit}
 \lim_{h\downarrow0}\frac{x^T(H'_h-H_h)x}{x^TH_hx}
   =\left(\frac43\right)^{1-\alpha}-1>0.
\end{equation}
For every $a>0$, let the exact response integral retain only shifts at
least $a$:
\[
 T_{h,a}=c_\alpha\int_a^\infty t^{\alpha-1}
 [F_t(L_{2,h},D_{2,h})-F_t(L_{1,h},D_{1,h})]\,dt.
\]
Then
\[
 0\le x^TT_{h,a}x
 \le\frac{8c_\alpha}{2-\alpha}\,h^2a^{\alpha-2},
 \qquad
 \lim_{h\downarrow0}\frac{x^TT_{h,a}x}{x^TH_hx}=0.
\]
In particular, an accuracy below the constant in \eqref{eq:limit}
cannot be obtained uniformly by discarding shifts below a positive
cutoff depending only on $n=4,\tau=1,\alpha$, and that accuracy.
\end{theorem}
\begin{proof}
The two cases $w=1,2$ of Lemma~\ref{lem:slow} give
\eqref{eq:limit}. For the cutoff bound, abbreviate $L=L_{1,h}$,
$L'=L_{2,h}$, and similarly for $D,D',K_t,K'_t$.
Since $x$ is constant across the edited edge,
\[
 Lx=L'x=z_h=(0,2h,-2h,0)^T,
 \qquad x^TLx=x^TL'x.
\]
The second expression for $F_t$ gives the exact identity
\[
 x^T(F'_t-F_t)x
 =\frac1t z_h^T(K_t^{-1}-(K'_t)^{-1})z_h.
\]
Now $K'_t\succeq K_t\succeq tI$, because $D\succeq I$.
The displayed quantity lies between zero and
$\|z_h\|^2/t^2=8h^2/t^2$. Integrating proves the bound.
Divide by $x^TH_hx\sim4h^\alpha$ to obtain the limit.
The omitted part is PSD by Theorem~\ref{thm:three}, so the relative
operator norm of its error is at least its Rayleigh quotient on $x$.
\end{proof}

\paragraph{Quantifiers and limits.}
This theorem rules out copying the existing positive-integral truncation
argument with an $n,\tau$-only low cutoff. It does not rule out adaptive
quadrature, an analytic treatment of low scales, additional cluster
corrections, a different representation, or a gap-free normalized
algorithm. It is not a computational lower bound. The example is weighted;
it does not prove the same obstruction for simple unweighted hosts.

\section{What remains a partitioning problem}
The generalized eigenvectors of $(H_\alpha,D)$ are exactly those of
$(L,D)$ and their eigenvalues are the $\alpha$ powers of the normalized
Laplacian eigenvalues. Thus fractional powering alone does not improve
Fiedler sweeps, prevent localization, or sharpen Cheeger rounding.
The new interface is a monotone common-kernel update and its determinant
budget. To yield a substantial partitioning advance, a future algorithm
must additionally resolve low-scale normalization and supply a new
rounding, certificate, or end-to-end computational guarantee.

\begin{thebibliography}{9}
\bibitem{KA} F.~Kubo and T.~Ando, Means of positive linear operators,
Math.~Ann.~246 (1980), 205--224.
\url{https://doi.org/10.1007/BF01371042}.
\bibitem{DB} F.~D\"orfler and F.~Bullo, Kron reduction of graphs with
applications to electrical networks, IEEE Trans.~Circuits Syst.~I
60 (2013), 150--163.
\url{https://arxiv.org/abs/1102.2950}.
\bibitem{BCKS} B.~Beckermann, A.~Cortinovis, D.~Kressner,
M.~Schweitzer, Low-rank updates of matrix functions II: Rational Krylov
methods, SIAM J.~Numer.~Anal.~59 (2021), 1325--1347.
\url{https://arxiv.org/abs/2008.11501}.
\end{thebibliography}
\end{document}
